\section{Rational centers and the cost of growing powers}
\label{sec:rational-centers}

For \(x\in[0,1)\), a positive integer \(a\) is a hit at \((h,s)\)
if its fractional part \(\{ax\}\) starts with \(V^s\) for some
length-\(h\) word \(V\). Write \(w=[V]_b\), \(u=\lfloor ax\rfloor\),
and \(P=aD v=Du+w\). The corresponding rational center is
\(v=P/(aD)\). The digit cylinder for \(ax\) is equivalent to
\begin{equation}\label{eq:signed-centre}
        -\frac{w}{M}\le aD(x-v)<\frac{D-w}{M}.
\end{equation}
Indeed its fractional interval is
\([Lw/M,(Lw+1)/M)\), and \(LD=M-1\). The first inequality
is strict when \(x\) is irrational. In particular
\(|x-v|<1/(aM)\) for irrational \(x\).
The \emph{reduced cofactor} of a center is
\[
 q=\operatorname{den}(Dv)=\frac{a}{\gcd(a,P)},
\]
where \(\operatorname{den}\) means the positive reduced denominator.
It divides both \(a\) and \(\operatorname{den}(v)\).

\subsection{Uniform lower bounds for the reduced cofactor}

\begin{lemma}\label{lem:late-three-forms}
For each fixed \(\xi>0\), outside a set of dyadic indices
satisfying \eqref{eq:bin-count}, all the corresponding \(h\)
satisfy
\begin{equation}\label{eq:late-three-forms}
 (a_\perp)^2\,\|a(b^h-1)b^m\alpha\|>B^{-\xi}
 \quad (0\le m\le T(h),\ 1\le a\le b^{T(h)}).
\end{equation}
\end{lemma}

\begin{proof}
Normalize \(0<\alpha<1\). For a counterexample use the nearest
integer \(p\) to \(a(B-1)b^m\alpha\), and put
\(e=a(B-1)b^m\alpha-p\). Divide
\((Ba b^m,a b^m,p)\) by \(g=\gcd(a b^m,p)\) to get
\((B k,k,p')\) with \(k=a b^m/g\), \(p'=p/g\).
The sup norm is \(B k\), and the real forms are the three
used in Theorem~\ref{thm:late-separation}. Exponents may be
taken as the exact logarithmic ratios for the values \(Bk,k,|e|/g\)
at the real place, and for the first two coordinates at the primes
of \(b\), with zero for the third at those primes.
The sum, multiplied by \(\log(Bk)\), is
\[
 \log\bigl((k_\perp)^2 |e|/g\bigr)
     \le -\xi\log B,
\]
because \(k_\perp\le a_\perp\). The third real exponent is
nonpositive, and the first equals \(1\). The log height lies in
\([h\log b,(h+2T(h))\log b]\), permitting the same choice of
\(\delta\) and the interval argument.

For any rational plane containing these vectors, its equation
\(uX_1+v_0X_2+w_0X_3=0\) with \(w_0=0\) fixes \(B\).
With \(w_0\ne0\), the points
\((B,e/(a b^m))\) lie on the line
\[
 y=(u/w_0+\alpha)x+(v_0/w_0-\alpha).
\]
Here the upper bound in \eqref{eq:secant-hyp} follows from
\(|e|\le B^{-\xi}\), and the lower bound follows by Liouville's
inequality with denominator \(a(B-1)b^m\le b^{h+2T(h)}\)
and division by \(a b^m\). The horizontal exponent in
\eqref{eq:secant-hyp} is \(\Delta=h\). This gives
\(O(1+\log F(J))\) indices per plane. Combined with the
height-interval estimate, this proves the lemma.
\end{proof}

\begin{theorem}\label{thm:reduced-cost}
Fix \(\xi>0\) and use the retained indices of
Theorem~\ref{thm:late-separation} and Lemma~\ref{lem:late-three-forms}.
Let \(x=\{b^m\alpha\}\) with \(0\le m\le T(h)\) and suppose
\(a\) is a hit with \(aD\le b^{T(h)}\).
For \(q=\operatorname{den}(Dv)\) the following inequalities hold:
\begin{equation}\label{eq:reduced-cost}
 \begin{aligned}
 q q_\perp &> \frac{a M}{D^2} B^{-\xi}
                      \ \ge a B^{s-2-\xi},\\
 q q_\perp^2 &> \frac{a M}{D} B^{-\xi}
                      \ \ge a B^{s-1-\xi}.
 \end{aligned}
\end{equation}
In particular
\begin{equation}\label{eq:qperp}
 a_\perp\ge q_\perp>
      B^{\max\{s-2-\xi,\ (s-1-\xi)/2\}} .
\end{equation}
For \(0<\xi<1\) the eventual period length of the canonical
base-\(b\) expansion of \(v\) exceeds
\(h\max\{s-2-\xi,(s-1-\xi)/2\}\).
\end{theorem}

\begin{proof}
Since \(q Dv\) is integral, \eqref{eq:signed-centre} gives
\[
 \|qD b^m\alpha\|=\|qD x\|< (q/a)D/M .
\]
Apply \eqref{eq:uniform-q} to \(qD b^m\), which is at most
\(b^{2T(h)}\) and has base-external part \(q_\perp D\).
Apply \eqref{eq:late-three-forms} with \(a\) replaced by \(q\).
This gives both inequalities in \eqref{eq:reduced-cost}.
Using \(q\le a\) gives \eqref{eq:qperp}.
If \(R=\operatorname{den}(v)\), then \(q_\perp\mid R_\perp\).
After the base-prime part is removed by advancing past the
preperiod, the eventual period \(t\) is the multiplicative order
of \(b\) modulo \(R_\perp>1\). Thus \(b^t-1\ge R_\perp\ge
q_\perp\), as claimed.
\end{proof}

For example \eqref{eq:reduced-cost} implies the explicit cancellation
bound
\[
 \gcd(a,P)<\sqrt{a}\,D\,B^{\xi/2}/\sqrt{M}.
\]
When \(a\le B^{s-\gamma}\) this is at most
\(B^{1-\gamma/2+\xi/2}\), uniformly in \(s\).
For an elementary digit comparison at the rational center, assume
\(0<\xi<1\). Then \(R_\perp>1\), and every \(b\)-adic endpoint
\(j/b^n\) is at distance at least \(1/(R b^n)\) from \(v\).
Consequently \(|x-v|<1/(aM)\) ensures that \(x\) and \(v\)
share the first \(n\) digits whenever \(b^n\le aM/R\).
In particular this gives \((s-1)h\) common leading digits, as
\(R\le aD\). For hits with \(a\le B^{s-\gamma}\) the
upper endpoint \(n=\lfloor\log_b(aM/R)\rfloor\) of this
guarantee is at most \((s+2+\xi-\gamma)h\) by
\(R\ge R_\perp\ge q_\perp > B^{s-2-\xi}\).
For \(s>6+3\xi\) this endpoint is shorter than two eventual
periods of \(v\), even before accounting for the preperiod.
The guarantee comes only from the cylinder bound. Actual
agreement could extend further.

\begin{corollary}\label{cor:leading-cost}
On the retained indices, let \(s=s(h)\to\infty\) with
\(2\le s\le F(j(h))/8\). Let \(a_*(h,s)\) be the least
positive integer multiplier giving a complete \(s\)-fold power
with root length \(h\) at some phase \(0\le m\le T(h)\).
The same conclusion holds if the minimum is taken only over the
phase window \eqref{eq:growing-phase-window}, with
\(a\equiv1\bmod b\) and least period \(h\). In either case
\[
       \frac{\log_b a_*(h,s)}{s h}\longrightarrow 1 .
\]
\end{corollary}

\begin{proof}
Theorem~\ref{thm:growing-power} gives
\(a_*(h,s)\le B^{s-\gamma}\) for all sufficiently large retained
\(h\), for any fixed \(0<\gamma<1/2\).
This upper bound ensures \(a_*D\le b^{T(h)}\).
The lower bound \eqref{eq:qperp} gives
\(\log_b a_*(h,s)>(s-2-\xi)h\). Divide by \(sh\).
\end{proof}

\subsection{When two multipliers have one center}

The next statement is elementary. Its arithmetic consequences
use the retained indices above. The square case with coprime
multipliers gives the threshold \(a+c\le B+2\).

\begin{proposition}\label{prop:common-centre}
Let \(x\) be irrational. If \(a,c>0\) are hits at the same
\((h,s)\) and
\begin{equation}\label{eq:centre-threshold}
       \frac{a+c}{\gcd(a,c)}\le L+1,
\end{equation}
then their rational centers coincide. If \(a\) is a hit of
center \(v\) and \(d\mid a\) with \(\operatorname{den}(Dv)\mid d\),
then \(d\) is also a hit with center \(v\).
In particular the common-center case transfers to \(\gcd(a,c)\)
and to \(\operatorname{den}(Dv)\).
\end{proposition}

\begin{proof}
In case of distinct centers, exchange \(a,c\) so that
\(v_a<v_c\), and divide \(a,c\) by \(g=\gcd(a,c)\) to get
\(a',c'\). Set \(r=a'P_c-c'P_a\), a positive integer.
From \eqref{eq:signed-centre} applied at both centers,
\[
       rM<c'(D-w_a)+a'w_c .
\]
The two sides are congruent modulo \(D\), since
\(P_a\equiv w_a\), \(P_c\equiv w_c\bmod D\), and \(M\equiv1\bmod D\).
Their gap is at least \(D\). Thus \(M+D\le(a'+c')D\),
and \(M+D=(L+1)D+1\), proving the first statement.

For the second, put \(t=a/d\). If \(x>v\), choose the residue
\(w_d\in[0,D-1]\) of \(dDv\) and take \(w_a\) from the hit.
The positive representative \(D-w_a\in[1,D]\) is at most
\(t(D-w_d)\) since these two positive integers are congruent
modulo \(D\). Thus the upper inequality for \(a\) in
\eqref{eq:signed-centre} gives the upper inequality for \(d\).
The lower inequality is automatic. If \(x<v\), use the
representative \(w_d\in[1,D]\) and the corresponding positive
\(w_a\in[1,D]\). Then \(w_a\le t w_d\), giving the lower
inequality for \(d\), while the upper is automatic. In either
case with \(u_d=(dDv-w_d)/D\in\Z\), the signed inequalities place
\(dx\) in the cylinder with integer part \(u_d\) and digit
label \(Lw_d\). This proves the transfer.
\end{proof}

\begin{proposition}[Sharpness]\label{prop:sharp-centre}
The bound \(L+1\) in \eqref{eq:centre-threshold} for equality of
centers is sharp for every \(h\) and \(s\).
For each fixed \(b\) and \(s\) the same bound for transfer to
the greatest common divisor is sharp for arbitrarily large \(h\).
\end{proposition}

\begin{proof}
For equality of centers, take \(a=1,\ c=L+1\) and irrational
\(x\in(L/(cM),1/M)\). The two root labels are \(0,1\) and their
centers are distinct.

For the transfer assertion choose \(1<a\mid D-1\) with
\(\gcd(a,L+2)=1\) and let \(c=L+2-a\).
This choice is possible for all large \(h\) for even \(b\),
with \(a=2\), because \(D-1\) is even and \(L+2\) is odd.
For odd \(b\), take \(a=B-2\), and choose \(h\) to be a
multiple of the multiplicative orders of \(b\) modulo the
primes \(p\mid 2^s+1\) with \(p\nmid b\). Indeed
\(L+2\equiv 2^s+1\bmod a\), and at each such prime \(B\equiv1\).
At a prime dividing both \(b\) and \(2^s+1\) one has \(B\equiv0\);
neither case allows \(p\mid a\).

Set \(w_a=0\), \(w_c=D-(D-1)/a\), so \(a w_c=1+D(a-1)\).
Choose \(1\le u<a\) with \(cu\equiv-1\bmod a\),
and \(v_0=(cu-a+1)/a\in\Z\). For the potential hits take
\(P_a=Du\), \(P_c=Dv_0+w_c\), so \(aP_c-cP_a=1\).
The proposed centers satisfy \(v_a=u/a<v_c\).
The right endpoint of the \(a\)-cylinder is greater than the
left endpoint of the \(c\)-cylinder precisely when
\(M<cD+a w_c\), and the latter sum is \(M+D\).
The right endpoint of the \(c\)-cylinder is at least \(v_c\),
so it is greater than the left endpoint \(v_a\) of the
\(a\)-cylinder. Their interiors therefore intersect.
Take irrational \(x\) in the intersection. Its distance to
\(u/a\) is less than \(1/(aM)\). Since \(a\mid D-1\) and
\(\gcd(cu,a)=1\), it follows that
\[
  \|Dx\|\ge 1/a-D/(aM)>D/M,
\]
using \(M>(a+1)D\). Thus \(1=\gcd(a,c)\) is not a hit.
Multiplying both multipliers by any positive \(g\) and replacing
\(x\) by \(x/g\) gives the same threshold for a prescribed
greatest common divisor \(g\).
\end{proof}

\begin{corollary}\label{cor:centre-obstructions}
Under the assumptions of Theorem~\ref{thm:reduced-cost} with
\(0<\xi<1\), two hits sharing a center cannot have a greatest
common divisor composed only of primes dividing \(b\).
Moreover if \(A D\le b^{T(h)}\) and \(2A\le L+1\), the number
of hits with \(1\le a\le A\) at one phase is at most
\[
 \min\!\left\{
      (A B^{2-s+\xi})^{1/2},\,
      (A^2 B^{1-s+\xi})^{1/3}
      \right\}.
\]
If this set is nonempty it has a unique common center and its
smallest multiplier is \(\operatorname{den}(Dv)\).
\end{corollary}

\begin{proof}
The reduced cofactor \(q\) divides the greatest common divisor
and has \(q_\perp>1\) by \eqref{eq:qperp}.
For the count, \(2A\le L+1\) implies \eqref{eq:centre-threshold}
for every pair.
All hits are multiples of \(q\), and \(q\) itself is a hit by
Proposition~\ref{prop:common-centre}. If \(a_{\max}\) is the
largest hit, the number is at most \(a_{\max}/q\).
In \eqref{eq:reduced-cost} one may substitute \(a_{\max}\)
and use \(q_\perp\le q\), so
\(q^2>a_{\max}B^{s-2-\xi}\) and
\(q^3>a_{\max}B^{s-1-\xi}\). The result follows using
\(a_{\max}\le A\).
\end{proof}

There is also a two-phase restriction. Suppose \(a\) hits at
\(m\), \(c\) hits at \(m+\Delta\), and both use root length \(h\)
and exponent \(s\). At the earlier tail \(x\) these are hits
of \(a\) and \(b^\Delta c\). Suppose these two hits of \(x\)
share the same center, and \(aD\le b^{T(h)}\).
Then \(q\mid\gcd(a,b^\Delta c)\) and the proof of
Theorem~\ref{thm:reduced-cost}, using the stronger precision
\(|x-v|<b^{-\Delta}/(cM)\), gives
\[
 b^\Delta c M B^{-\xi}
   < \min\{D^2 q q_\perp,\ D q q_\perp^2\}.
\]
Thus at retained \(h\) even a later phase sharing an earlier
center needs a substantial base-external common factor.
These restrictions concern this specific transfer through a
rational center. The existence of sufficiently late powers or
palindromes in the original expansion remains open.

